\section{The common positive quartic and its literal index}
\label{sec:parent}

We use the unshifted positive quartic of package 30 as the coding source.
This choice keeps the inherited coefficient index unchanged as a
symbolic recipe.  The cheaper 60-operation system later in this report
is an equivalent presentation; it is not silently substituted into the
index compiler.

Let $\beta,K_0,M_C,M_F$ denote the complete tuple of
\eqref{eq:new-four-ports}; $K_0$ is that equation's port $K$ and is distinct
from subsequent polynomial expressions also traditionally called $K$.
The 31 positive coordinates, in their fixed order, are
\[
 (J,F,\alpha,z,f,h,i,o,j,s,w,\tau,\eta,\zeta,y,\gamma,
 N,E,P,c,a,X,D,C_2,L,F_2,V,g,I,B_{\mathrm{aux}},\Delta).
\]
Put
\[
 q=\beta J+1,\quad k=\eta+\zeta,\quad
 H=4a+3,\quad C=q-F-\alpha.
\]
The parent is the sum of the squares of the residuals of these 22
equations, with ``left minus right'' as the residual convention:
\begin{alignat}{2}
 L&=q^2-1, &\qquad N&=q(L+1),\notag\\
 X&=wq, & a&=E+sN,\notag\\
 E&=X(a-E), & c&=k(a-E)+\eta,\notag\\
 P&=E(c-\eta)+k, & \tau^2&=P(P-k)+1,\notag\\
 D&=X+ac+\gamma H, & \Delta&=a^2+H,\notag\\
 C_2&=c^2, & D^2&=\Delta C_2+1,\notag\\
 k&=jc-V+1+hE, & F_2&=f^2,\notag\\
 V&=of-c, & g&=y^2-V^2,\notag\\
 I&=iC_2, & B_{\mathrm{aux}}&=I^2,\notag\\
 B_{\mathrm{aux}}&=\Delta(F_2-1), & y^2-B_{\mathrm{aux}}g&=1,
 \label{eq:parent-quadratic}
\end{alignat}
together with the two transport equations
\begin{align}
 jc-V&=(q-F)(N-q)-CL+(M_C+qM_F)J,\label{eq:parent-mask}\\
 (K_0+w)C+q-F-z(q-1)&=1.\label{eq:parent-transport}
\end{align}
The order in the literal JSON, rather than the grouping on this page,
is authoritative for the coefficient indices.  There are no products of
more than two coordinates in a residual after the displayed abbreviations
are substituted: every residual has total degree at most two.

Write this sum of squares as $Q_{31}$ and let
\[
 P_\star(x_0,x_1,\ldots,x_{31})
   =Q_{31}(x_1+1,\ldots,x_{31}+1).
\]
The dummy input $x_0$ does not occur.  The coefficient of $\tau^4$ in
$Q_{31}$ is 1, so the degree is exactly four.  At the all-zero
nonnegative input the residual $E-X(a-E)$ is 1, hence
$P_\star(a,0,\ldots,0)\geq1$ for every integer $a$.  Thus Sun's source
construction applies directly; there is no need to append a variable or
square this quartic again.

\begin{proposition}\label{prop:parent}
With the complete tuple \eqref{eq:new-four-ports}, ZFC proves
\[
 (\exists\boldsymbol{n}\in\N^{31}\ P_\star(0,\boldsymbol{n})=0)
 \ \longleftrightarrow\ 
 (\text{there is a finite ZFC refutation}).
\]
\end{proposition}
\begin{proof}
Translation by 1 is a bijection from $\N$ to $\Z_{>0}$.  A sum of integer
squares is zero exactly when every residual is zero.  The package-30
parent proof identifies these positive solutions with the compatible
native arithmetic skeleton.  Section~\ref{sec:new-coefficients} supplies
that skeleton with a complete tuple preserving its soundness,
completeness, slack and parity interface.  The native compiler and the
K215 endpoint then give the assertion.  The strong auxiliary Pell index
must be converted to the ordinary index as in the inherited parent
proof; the tuple replacement does not change that conversion.
\end{proof}

\subsection{Weighted coefficients with no unspecified index}

Put $\nu=31$ and $L_\star=5^{32}=23283064365386962890625$.  Write
\[
 P_\star=\sum_{|\boldsymbol{a}|\le4}p_{\boldsymbol a}
                x_0^{a_0}\cdots x_{31}^{a_{31}},\qquad
 c_{\boldsymbol a}=(4-|\boldsymbol a|)!
              \prod_{j=0}^{31}a_j!\ p_{\boldsymbol a}.
\]
Define $e(\boldsymbol a)=\sum_{j=0}^{31}a_j5^j$.
All $58,905=\binom{36}{4}$ positions of total weight at most four are
distinct, and $e(\boldsymbol a)<L_\star$.  There are 402 nonzero weighted
coefficient records when the four ports are treated as fixed symbolic
parameters.  The file \texttt{signed16/data/constants.json} gives all of
them as explicit finite recipes.  The verifier expands the normalized
parent symbolically and checks every record, including factorial weights.

For the new tuple, the explicitly materialized integer
\[
 z_\star=2^{4269855}
\]
satisfies
\begin{equation}\label{eq:coefficient-bound}
 |c_{\boldsymbol a}|+1<z_\star
 \quad\text{for every }|\boldsymbol a|\le4.
\end{equation}
The check uses exact integer arithmetic on the actual port payloads,
not a floating-point estimate.  This also bounds every
$|p_{\boldsymbol a}|$ by $z_\star$, since its factorial multiplier is a
positive integer.

For the Jones route define the fixed polynomial and fixed integers
\begin{align*}
 e_\star(T)&=\sum_{|\boldsymbol a|\le4}
                  (z_\star+c_{\boldsymbol a})T^{L_\star-e(\boldsymbol a)},\\
 y_\star&=e_\star(2z_\star),\qquad
 \kappa=(2z_\star)^{L_\star+1}
        =2^{99415332078933715820316769856}.
\end{align*}
Every digit lies strictly between 0 and $2z_\star$ and its position is
distinct, so $0<y_\star<\kappa$.  These are bounded definitions of actual
integers.  The huge $y_\star$ need not be a leaf of a circuit, and the
power-of-two recipe for $\kappa$ is explicit even though expanding its
binary string is not practical.

For the ten-variable route we instead use the sparse polynomial
\[
 D_\star(T)=\sum_{|\boldsymbol a|\le4}
                   c_{\boldsymbol a}T^{L_\star-e(\boldsymbol a)}.
\]
It has exactly the same weighted coefficient data, without the all-position
$z_\star$ padding.  This is one reason to preserve the literal parent and
its coordinate order across the two independent reductions.
